\section{Introduction}\label{sec:intro}

OpenAI's openai/math repository publishes mathematical manuscripts together with Lean formalizations and Comparator challenges; we use it at commits adc7f12 and fd4aeeb \cite{openai-math-adc7f12,openai-math-fd4aeeb}.
The manuscript ``Uniform exclusion of Landau--Siegel zeros'' (OpenAI, October 1, 2026) \cite{openai-siegel-2026} states as its Theorem 1 that there is an absolute constant c > 0 such that every real zero beta in (0,1) of every primitive nonprincipal real Dirichlet L-function of conductor q >= 3 satisfies (1 - beta) log q >= c.
Comparator's README \cite{comparator} states that for each listed theorem it guarantees three things: the theorem proves the statement given in the challenge, uses no more axioms than permitted, and is accepted by the Lean kernel.
Comparator can also pass the solution export to an external kernel; we used nanoda \cite{nanoda-lib}, at the official build 3a24072, as the second kernel.

Acceptance shows that a Lean proof of the Lean statement exists, using only the permitted axioms.
It does not show that the Lean proof follows the paper's argument: in the case studied here, both accepted OpenAI proofs obtain the paper's spanning step by other means (Section~\ref{sec:gap}).
Bastounis, Circelli and Hansen \cite{bastounis2026} argue that a Lean check of an AI autoformalisation need not certify the natural-language proof, and show for OpenAI's Navier–Stokes formalisation that the Lean proof does not correspond to the natural-language proof.
Their Section 1 states that they make no claims about the correctness of OpenAI's natural-language proof.

We find that the challenge has two checked theorems, with equivalent statements, and that OpenAI proves them by two routes.
OpenAI's catalogued main result, exists\_absolute\_real\_zero\_gap, gets its spanning step from a local B\'ezout statement, through a zero test with Corollary 4's bounds, and does not use OpenAI's Corollary 4 theorem.
The second theorem, dirichletRealZeroBound\_proof, proves Corollary 4 as actual\_biquadratic\_rectangle\_span by a multiplicity estimate, uniform\_rectangular\_multiplicity.
So neither route goes through the paper's proof of Lemma 3.

Our contributions are these.
\begin{itemize}
\item Statement faithfulness of the four family-003 challenges, by reading the Lean statements against the papers' texts, with Lean lemmas showing the SiegelZeros challenge is not vacuous (Section~\ref{sec:faithfulness}, with the reasoning behind each reading in Tables~\ref{tab:points-SiegelZeros}--\ref{tab:points-HeckeSevenEighths}).
\item Lean proofs of the paper's Lemma 3 and Corollary 4, with statements that match the paper's and only the standard axioms (Section~\ref{sec:lemma3}).
\item A reroute of dirichletRealZeroBound\_proof through that Corollary 4, accepted by Comparator with both kernels (Section~\ref{sec:reroute}).
\item A replay of the openai/math corpus as context, compared with an independent replay (Section~\ref{sec:corpus}).
\end{itemize}

We do not claim the following.
\begin{itemize}
\item Nothing about the paper's analytic arguments, nor about either OpenAI route beyond its spanning step.
\item No specialist review of the mathematics: no specialist in analytic number theory or transcendence has reviewed it.
\item Statement faithfulness rests on reading, by the author, with the reasoning recorded row by row.
\item The reroute changes one of the two checked theorems; OpenAI's main result is untouched.
\item The separate copy of this code in namespace OAI.Erdos970, loaded by 3 configurations, is outside our scope.
\item For an explicit value of the constant we point to the note docs/effectivity/note-v2.pdf in the reviewer-packet-v3 archive \cite{zenodo-23523697} and to \cite{arxiv-2610-15483}, and make no claim about it here.
\end{itemize}

Rerouting OpenAI's main result through the paper's Lemma 3 would require a bridge from Lemma3.rectangle to ActualCutoffSpanning; we leave this open.

The paper is organized as follows.
Section~\ref{sec:background} describes Comparator, nanoda, the challenge and the paper's lemmas.
Section~\ref{sec:faithfulness} reads the four challenge statements against the papers.
Section~\ref{sec:gap} describes OpenAI's two formal routes.
Section~\ref{sec:lemma3} gives the Lean proofs of Lemma 3 and Corollary 4.
Section~\ref{sec:reroute} reports the reroute.
Section~\ref{sec:corpus} places it in the corpus replay.
Section~\ref{sec:threats} lists threats to validity, Section~\ref{sec:artifacts} the artifacts, and Section~\ref{sec:ai} the use of AI tools.
